% Dot batch 84; source SHA-256: fd7c0f65cd58e6e6c531c7e039a2e7eefce14a96b48340669cd1fa7e69c91439
% Archived source: dot_updates_20261124/batches/batch084/EGA第五卷_任意底上的单态性与无穷小差.tex
% Original body lines: 28--86; integration map: integration/late_integration.json
\begingroup
\def\newscope{V-L}
\def\Spec{\operatorname{Spec}}
\def\Aut{\operatorname{Aut}}
\def\Div{\operatorname{Div}}
\def\OO{\mathcal O}
\def\PP{\mathbf P}
\def\m{\mathfrak m}
\def\Hom{\operatorname{Hom}}
\noindent\textbf{命题4.3的证明：任意底。}
现在概述一般情形，即 \(S\) 任意。换底之后，a) 只须证明：\(P^\vee\) 在 \(S\) 上定义了 \(X\) 上同一个除子的两条截面必相等；b) 只须证明：满足 \(D\circ u=D\circ v\) 的 \(P^\vee\) 的两个自同构必相等。
可以假定 \(S\) 仿射。在 b) 中，没有专门假定 \(g\) 有限表示，但 \(g\) 满射。利用 \(g\) 是开映射，可以局部归结到 \(X\) 也仿射，因而在 \(S\) 上有限表示的情形。再用惯用的逼近方法，归结到 \(S\) 诺特的情形。\cmt{L18-084-01}{法143的“可以假定 \(S\) 仿射”为手写补入。b) 归约时，取 \(X\) 的仿射开，其像覆盖底；围绕底点选入该像内的主开，并限制该仿射开到这个主开。不能任取一个仿射开就断言它仍满射到底。原稿把诺特逼近留作惯用步骤；本版核查报告 L18-084 明确比较的有限数据。}

设底概形 \(S\) 诺特，且有 \(S\) 上的函子态射 \(h:F\to G\)；\(F,G\) 都是从 \((\mathrm{Sch}/S)^{\mathrm{op}}\) 到集合范畴的函子。有很一般的判据（将在第五章阐明）：在关于 \(F,G\) 的一些简单假设下，如果每个 \(s\in S\) 上的纤维态射 \(F_s\to G_s\) 都是单态射，则 \(F\to G\) 是单态射。这里记 \(F_s=F\times_S\Spec k(s)\)，\(G_s\) 同理。举例说，当两者都由 \(S\) 上局部有限型的概形表出时，这些假设成立；但在我们这里，先验有表出的只有第一个函子。现就所关心的两个特例概述第五章的论证\nobreak\cmt{L18-084-02}{原稿说的是纤维函子的单态射，也就是对纤维上的所有检验概形单射，而不仅是剩余域点的单射。英108漏写了法143例子中的“局部”。这里的第五章是手稿中预告的章节。}。

我们有有限型 \(S\)-概形 \(F\) 的两个截面 \(u,v\)（分别取 \(F=P^\vee\) 或射影群 \(F=\mathrm{GP}(E^\vee)\)），要证明二者相等。只须对每个 \(s\in S\) 和每个 \(n\geq0\)，在换底
\[
 \Spec\bigl(\OO_{S,s}/\m_s^{n+1}\bigr)\longrightarrow S
\]
后证明两截面相等。\cmt{L18-084-03}{法143与英108一笔带过，此处由闭对角线形成等化子，再用克鲁尔交定理；见本版核查报告 L18-084。}
于是归结到阿廷局部底 \(S=\Spec A\)。设 \(\m^{n+1}=0\)，按 \(n\) 归纳。可假定 \(u,v\) 模 \(\m^n\) 相等，于是它们相差一个元素
\[
 \delta\in\Hom_k\bigl(u_0^*\Omega^1_{F_0/k},V\bigr).
\]
这里 \(k=A/\m\) 是剩余域，\(F_0=F\times_A\Spec k\) 是剩余域纤维，\(V=\m^n=\m^n/\m^{n+1}\)。利用 \(h(u)=h(v)\)，只须证明 \(\delta=0\)。\cmt{L18-084-04}{法144、英109称 \(F_0\) 为“约化纤维”。公式只对底取剩余域，并不把所得纤维再取约化，否则可能丢失切向量。归纳的 \(n=0\) 步直接用纤维单态射；对 \(n\ge1\)，\(V^2=0\) 且 \(\m V=0\)。}

验证的一般原则如下。先把 \(h(u),h(v)\) 看成模 \(\m^n\) 一致的两个点。它们的差可写成
\[
 \delta'\in\Hom_k\bigl(w_0^*\Omega^1_{G_0/k},V\bigr),
 \qquad w_0=h_0(u_0)=h_0(v_0),\quad G_0=G\times_A\Spec k.
\]
这个元素恰由 \(\delta\) 与 \(h_0:F_0\to G_0\) 所诱导的自然映射
\[
 h_0^*:w_0^*\Omega^1_{G_0/k}\longrightarrow u_0^*\Omega^1_{F_0/k}
\]
复合得到。因为 \(h(u)=h(v)\)，\(\delta'=0\)，即 \(\delta\) 与 \(h_0^*\) 的复合为零。因此只要 \(h_0^*\) 满射，便有 \(\delta=0\)。\(h_0\) 是单态射，特别是在 \(k\) 的对偶数环值点上诱导单射，所以 \(h_0^*\) 的对偶单射，\(h_0^*\) 也就满射。\cmt{L18-084-05}{法144的中间一句漏写星号，字面写成“\(h_0\) 满射”；所需的是前一行的余切映射 \(h_0^*\) 满射，而不是函子态射满射。英109同样有此漏写。}

这一论证可原样用于可表出的 \(G\)，但所关心的情形不属于此类。尽管如此，可定义一个 \(k\)-向量空间 \(T_{w_0}G_0\)，充当 \(w_0^*\Omega^1_{G_0/k}\) 的对偶所扮演的角色，即 \(G_0\) 在 \(w_0\) 的切空间。把两个模 \(\m^n\) 一致的 \(G\) 的点的偏差，表示为 \(T_{w_0}G_0\otimes_kV\) 的元素。这是直接的计算，已包含在手稿拟援引的无穷小扩张部分。\cmt{L18-084-06}{法144手写补入“切空间”及张量积。此处节号有涂改，英109暂作§26并带问号；下一页作者说明所引“小计算”的节号按英110读作§25。两处均属手稿内部预告，不改作正式出版章节。}

在 a) 中，\(G\) 是 \(\Div_{X/S}\)。于是 \(w_0\) 对应于 \(X_0=X\times_A\Spec k\) 上的 Cartier 除子 \(D_0\)，所取的切空间是
\[
 T_{w_0}G_0=H^0(D_0,\mathcal N_{D_0/X_0}),
 \qquad
 \mathcal N_{D_0/X_0}\simeq\OO_{X_0}(D_0)|_{D_0}.
\]
这里 \(\mathcal N_{D_0/X_0}\) 是 \(D_0\) 在 \(X_0\) 中的法丛。

在 b) 中，可以把 \(X\) 限制到共同基点集的补开，因而设 \(D\) 无基点。更便于用到 \(P\) 的态射来理解此情形（见4.1）。这时 \(G\) 是相对态射函子 \(\underline{\Hom}_S(X,P)\)，而 \(T_{w_0}G_0\) 是\cmt{L18-084-07}{法145、英109把 Hom 的靶写成 \(\OO_{P_0}\)。左项已拉回到 \(X_0\)，差分是到 \(X_0\) 上结构层的导子，故正文改为 \(\OO_{X_0}\)。取基点补开并不要求一般底上的仿射开本身全局满射。}
\[
 \Hom_{\OO_{X_0}}\bigl(f_0^*\Omega^1_{P_0/k},\OO_{X_0}\bigr).
\]


在两个情形中，都有自然线性映射
\[
 T_{u_0}F_0\otimes_kV\longrightarrow T_{w_0}G_0\otimes_kV,
 \qquad T_{u_0}F_0=\bigl(u_0^*\Omega^1_{F_0/k}\bigr)^\vee .
\]
它表示应用 \(h\) 后，偏差 \(\delta\) 变为对应的 \(\delta'\)。
该映射是单射，因为 \(T_{u_0}F_0\to T_{w_0}G_0\) 是单射；后者来自 \(h_0:F_0\to G_0\) 是单态射。

\medskip\noindent\textbf{作者的写作说明。}
实际撰写时，证明的最后一部分似乎难以不引用第25节的几个小计算；没有必要在此特例中重作这些计算。应指出，这不会构成循环论证，因为第25节及其中发展的计算只依赖第16节的微分演算结果；此外，在第四章中将不再使用4.3，至多可能在接下来的两小节使用。\cmt{L18-084-08}{这一段在法145左侧有斜线，似为编辑标记。英110编辑者猜测这部分原拟作 EGA 第四章的片段；这只是编排猜测，不属于命题。}

4.3 a) 的意义在于，在所述条件下，参数射影丛可解释为 \(\Div_{X/S}\) 的一个子函子。这样便可以内在地解释 \(X/S\) 上的线性系的同构类（同构允许改变参数射影丛）：它是 \(\Div_{X/S}\) 的一个子函子 \(P^\vee\)，由射影丛表出，而且由其典范包含所定义的除子族在第3小节的意义下为线性。这实质上是经典观点：除子线性系就是满足一定条件的一组除子（比较4.4）。\cmt{L18-084-09}{法145恢复两处 \(\Div_{X/S}\)。文中的“改变参数丛”指在同构携带除子族的意义下，并不忘记两个参数丛间的相容性。法稿末尾节号有覆写，英110取3；这里的“线性”沿用前文定义。}

另一方面，在 \(g\) 满射时，4.3 b) 也就是说，如果分别以两个射影丛 \(P^\vee,(P^\vee)'\) 参数化的 \(X/S\) 上的两个除子线性系同构，则存在唯一的 \(P^\vee\to(P^\vee)'\) 的同构，与 \(D,D'\) 相容。因此，可以说 \(X/S\) 上线性系的一个同构类确定其参数射影丛，精确到唯一同构。在技术上，等第五章的下降理论可用时，就能据此对除子线性系作忠实平坦下降；不过参数丛还须容许“扭曲射影丛”，这将在后续小节处理。\cmt{L18-084-10}{法146手写补入“当 \(g\) 满射时”，英110保留。原文在这里预告下降，并未证明全面的可表出性或下降定理。}
\par\endgroup
